\subsection{Order Statistics, Extrema, and Moment Methods}
% Source integration: attached screenshots from Chapter 4, "Expected Values",

Order statistics describe the sorted values of an independent sample. They turn joint density calculations into formulas for minima, maxima, medians, ranges, and distribution-free tolerance intervals.

\begin{definition} \label{def:order-statistics} \hlt{Order Statistics}\\
For real-valued $X_1,\ldots,X_n$,
\begin{enumerate}[label=(\roman*)]
\setlength{\itemsep}{0pt}
\item The \hlt{order statistics} are the pointwise rearrangement $X_{(1)}\leq X_{(2)}\leq\cdots\leq X_{(n)}$ of $X_1,\ldots,X_n$.
\item $X_{(1)}=\min_{1\leq i\leq n}X_i$, $X_{(n)}=\max_{1\leq i\leq n}X_i$, and if $n=2m+1$, then $X_{(m+1)}$ is the sample median.
\end{enumerate}
\end{definition}

\begin{proposition} \label{prop:min-max-cdfs-and-densities} \hlt{Minimum and Maximum of an IID Continuous Sample}\\
Let $n\geq1$ and $X_1,\ldots,X_n\overset{\mathrm{iid}}{\sim}F$, where $F$ is differentiable with density $f=F'$.\\
Set $V=X_{(1)}=\min_iX_i$ and $U=X_{(n)}=\max_iX_i$; then for $u,v\in\mathbb{R}$, $F_U(u)=F(u)^n$, $f_U(u)=n f(u)F(u)^{n-1}$, $F_V(v)=1-(1-F(v))^n$, and $f_V(v)=n f(v)(1-F(v))^{n-1}$.
\begin{proof}
Since $\{U\leq u\}=\bigcap_{i=1}^n\{X_i\leq u\}$, independence gives $F_U(u)=\prod_{i=1}^nF(u)=F(u)^n$.\\
Differentiation gives $f_U(u)=n f(u)F(u)^{n-1}$.\\
Since $\{V>v\}=\bigcap_{i=1}^n\{X_i>v\}$, independence gives $1-F_V(v)=\prod_{i=1}^n(1-F(v))=(1-F(v))^n$.\\
Differentiation of $F_V(v)=1-(1-F(v))^n$ gives $f_V(v)=n f(v)(1-F(v))^{n-1}$.\\
\hlt{Structural Commentary}: Extremes convert order events into products over independent coordinates.\\
\hlt{Key Insight}: Maxima use all observations to the left; minima use all observations to the right.
\end{proof}
\end{proposition}

\begin{example} \label{ex:series-parallel-exponential-lifetimes} \hlt{Series and Parallel Exponential Lifetimes}\\
Let $T_1,\ldots,T_n\overset{\mathrm{ind}}{\sim}\operatorname{Exp}(\lambda)$ be component lifetimes.
\begin{enumerate}[label=(\roman*)]
\setlength{\itemsep}{0pt}
\item
For a series system, all components must survive, so system failure occurs at the first component failure: $V=\min_i T_i=T_{(1)}$. Thus $F_V(v)=1-e^{-n\lambda v}$ and $f_V(v)=n\lambda e^{-n\lambda v}$ for $v\geq0$; hence $V\sim\operatorname{Exp}(n\lambda)$.
\item
For a parallel system, one surviving component suffices, so system failure occurs only after all components fail: $U=\max_i T_i=T_{(n)}$. Thus $F_U(u)=(1-e^{-\lambda u})^n$ and $f_U(u)=n\lambda e^{-\lambda u}(1-e^{-\lambda u})^{n-1}$ for $u\geq0$.
\end{enumerate}
\end{example}

\begin{theorem} \label{thm:kth-order-statistic-density} \hlt{Density of the kth Order Statistic}\\
Let $X_1,\ldots,X_n$ be iid real-valued random variables with differentiable CDF $F$ and density $f=F'$.\\
For $1\leq k\leq n$ and $x\in\mathbb{R}$, $f_{X_{(k)}}(x)=\frac{n!}{(k-1)!(n-k)!}f(x)F(x)^{k-1}(1-F(x))^{n-k}$.
\begin{proof} Set $N_x=\sum_{i=1}^n\mathbf{1}_{\{X_i\leq x\}}$.\\
Then $N_x\sim\operatorname{Bin}(n,F(x))$ and $\{X_{(k)}\leq x\}=\{N_x\geq k\}$, so $F_{X_{(k)}}(x)=\sum_{j=k}^n\binom{n}{j}F(x)^j(1-F(x))^{n-j}$.\\
With $p=F(x)$, $\frac{d}{dp}\sum_{j=k}^n\binom{n}{j}p^j(1-p)^{n-j}=\frac{n!}{(k-1)!(n-k)!}p^{k-1}(1-p)^{n-k}$, and the chain rule gives stated density.\\
\hlt{Structural Commentary}: The multinomial count of left, middle, and right observations supplies the coefficient.\\
\hlt{Key Insight}: An order statistic is located by counting how many sample points fall on each side.
\end{proof}
\end{theorem}

\begin{corollary} \label{cor:uniform-order-statistics-are-beta} \hlt{Uniform Order Statistics Are Beta}\\
If $X_1,\ldots,X_n\overset{\mathrm{iid}}{\sim}\operatorname{Uniform}(0,1)$, then for $1\leq k\leq n$, $X_{(k)}\sim\operatorname{Beta}(k,n-k+1)$\\
with density $f_{X_{(k)}}(x)=\frac{n!}{(k-1)!(n-k)!}x^{k-1}(1-x)^{n-k}\mathbf{1}_{(0,1)}(x)$.
\begin{proof}
On $(0,1)$, $F(x)=x$ and $f(x)=1$, so Theorem~\ref{thm:kth-order-statistic-density} gives $f_{X_{(k)}}(x)=\frac{n!}{(k-1)!(n-k)!}x^{k-1}(1-x)^{n-k}$.\\
Also $B(k,n-k+1)=\frac{(k-1)!(n-k)!}{n!}$, hence this density is the $\operatorname{Beta}(k,n-k+1)$ density.\\
\hlt{Structural Commentary}: Uniform order statistics reduce the general density to the beta normalizing constant.\\
\hlt{Key Insight}: Counting positions in a uniform sample produces beta powers of left and right mass.
\end{proof}
\end{corollary}

\begin{proposition} \label{prop:min-max-joint-density-and-uniform-range} \hlt{Joint Density of Minimum and Maximum}\\
Let $n\geq2$, let $X_1,\ldots,X_n\overset{\mathrm{iid}}{\sim}F$ with differentiable CDF $F$, density $f=F'$; $V=X_{(1)}$, $U=X_{(n)}$, $R=U-V$.
\begin{enumerate}[label=(\roman*)]
\setlength{\itemsep}{0pt}
\item For $v<u$, $f_{V,U}(v,u)=n(n-1)f(v)f(u)(F(u)-F(v))^{n-2}$, and $f_{V,U}(v,u)=0$ for $v\geq u$.
\item If $F(t)=t$ on $[0,1]$, then $f_{V,U}(v,u)=n(n-1)(u-v)^{n-2}\mathbf{1}_{\{0<v<u<1\}}$.
\item In the uniform case, $f_R(r)=n(n-1)r^{n-2}(1-r)\mathbf{1}_{(0,1)}(r)$, with $F_R(r)=0$ for $r<0$,\\
$F_R(r)=nr^{n-1}-(n-1)r^n$ for $0\leq r\leq1$, and $F_R(r)=1$ for $r\geq1$.
\end{enumerate}
\begin{proof}
\begin{enumerate}[label=(\roman*)]
\setlength{\itemsep}{0pt}
\item For $v<u$, $G(v,u)=\mathbb{P}(V\geq v,U\leq u)=\mathbb{P}(v\leq X_i\leq u,\ 1\leq i\leq n)=(F(u)-F(v))^n$,\\
so $f_{V,U}(v,u)=-\partial_v\partial_uG(v,u)=n(n-1)f(v)f(u)(F(u)-F(v))^{n-2}$.
\item If $F(t)=t$ and $f(t)=1$ on $(0,1)$, by (i), $f_{V,U}(v,u)=n(n-1)(u-v)^{n-2}$ on $0<v<u<1$, else $0$.
\item For $0<r<1$, $f_R(r)=\int_0^{1-r}f_{V,U}(v,v+r)\,dv=n(n-1)r^{n-2}(1-r)$.\\
Integration gives $F_R(r)=\int_0^r f_R(s)\,ds=nr^{n-1}-(n-1)r^n$; outside $[0,1]$, the support gives $F_R(r)\in\{0,1\}$.
\end{enumerate}
\hlt{Structural Commentary}: Mixed derivatives select the two boundary observations.\\
\hlt{Key Insight}: Fixed extremes leave only interior mass.
\end{proof}
\end{proposition}

\begin{example} \label{ex:distribution-free-tolerance-interval-from-range} \hlt{Distribution-Free Tolerance Interval}\\
Let $X_1,\ldots,X_n$ be iid with continuous CDF $F$, and set $Q=F(X_{(n)})-F(X_{(1)})$.\\
By the probability integral transform, $U_i=F(X_i)\overset{\mathrm{iid}}{\sim}\operatorname{Uniform}(0,1)$, so $Q\stackrel{d}{=}U_{(n)}-U_{(1)}$.\\
Therefore, for $0\leq\alpha\leq1$, $\mathbb{P}(Q>\alpha)=1-F_R(\alpha)=1-n\alpha^{n-1}+(n-1)\alpha^n$. For $n=100$ and $\alpha=.95$, $\mathbb{P}(Q>.95)\approx.96$: $[X_{(1)},X_{(n)}]$ covers at least $.95$ of the $F$-mass with probability about $.96$.
\end{example}

The expectation calculus developed above also supports a useful collection of moment methods: identifying whether the mean exists, replacing difficult distributions of functions by direct expectation, controlling spread by covariance, and using generating functions or Taylor expansion when direct integration is cumbersome.

\begin{example} \label{ex:expectation-existence-and-classical-means} \hlt{Existence and First Distributional Means}
\begin{enumerate}[label=(\roman*)]
\setlength{\itemsep}{0pt}
\item If $X\sim\operatorname{Geom}(p)$ on $\{1,2,\ldots\}$ and $q=1-p$, then $\mathbb{E}X=p\sum_{k\geq1}kq^{k-1}=p\frac{d}{dq}\frac{q}{1-q}=1/p$.
\item If $X\sim\operatorname{Poisson}(\lambda)$, then $\mathbb{E}X=e^{-\lambda}\sum_{k\geq1}k\lambda^k/k!=\lambda e^{-\lambda}\sum_{j\geq0}\lambda^j/j!=\lambda$.
\item If $X\sim\operatorname{Cauchy}(0,1)$, then $\mathbb{E}|X|=\frac{2}{\pi}\int_0^\infty x(1+x^2)^{-1}\,dx=\infty$, so $X$ is not integrable and $\mathbb{E}X$ is undefined.
\item If $\mathbb{P}(Y=2^k)=2^{-(k+1)}$ for $k\geq0$, then $\mathbb{E}Y=\sum_{k\geq0}2^{-1}=\infty$; truncation is essential.
\end{enumerate}
\hlt{Key Insight}: Expectation is a finite absolute average, not merely a formal symmetry or a plausible center.
\end{example}

\begin{method} \label{meth:indicator-sum-expectation-without-independence} \hlt{Indicator Sums Without Independence}
\begin{enumerate}[label=(\roman*)]
\setlength{\itemsep}{0pt}
\item For a finite index set $I$ and events $(A_i)_{i\in I}$, if $W=\sum_{i\in I}\mathbf{1}_{A_i}$, then $\mathbb{E}W=\sum_{i\in I}\mathbb{P}(A_i)$.
\item The identity uses only linearity of expectation; the events $A_i$ need not be independent.
\end{enumerate}
\end{method}

\begin{example} \label{ex:dna-word-count-by-linearity} \hlt{Counting Word Occurrences}\\
Let $S_1,\ldots,S_N\overset{\mathrm{iid}}{\sim}\operatorname{Uniform}\{A,C,G,T\}$ and fix a word $w\in\{A,C,G,T\}^q$.\\
For $1\leq m\leq N-q+1$, set $I_m=\mathbf{1}_{\{(S_m,\ldots,S_{m+q-1})=w\}}$ and $W=\sum_{m=1}^{N-q+1}I_m$.\\
Then $\mathbb{P}(I_m=1)=4^{-q}$ and $\mathbb{E}W=(N-q+1)4^{-q}$.\\
Overlaps may create dependence; linearity does not require independence.
\end{example}

\begin{example} \label{ex:coupon-collector-expectation-harmonic} \hlt{Coupon Collection by Waiting Times}\\
Draw uniformly from $n$ coupon types, and let $X_r$ be the waiting time for the $r$th new type. Then $X_r\sim\operatorname{Geom}((n-r+1)/n)$ on $\{1,2,\ldots\}$, so $\mathbb{E}X_r=n/(n-r+1)$.\\
For $T=\sum_{r=1}^nX_r$, linearity gives $\mathbb{E}T=\sum_{r=1}^n n/(n-r+1)=nH_n$, where $H_n=\sum_{j=1}^n1/j\sim\log n+\gamma$.
\end{example}

\begin{theorem} \label{thm:correlation-bound-and-equality} \hlt{Correlation Bound and Equality Case}\\
Let $X,Y\in L^2$ have means $\mu_X,\mu_Y$ and standard deviations $0<\sigma_X,\sigma_Y<\infty$, set $\rho=\operatorname{Cov}(X,Y)/(\sigma_X\sigma_Y)$, and define $Z_X=(X-\mu_X)/\sigma_X$, $Z_Y=(Y-\mu_Y)/\sigma_Y$.
\begin{enumerate}[label=(\roman*)]
\setlength{\itemsep}{0pt}
\item $\mathbb{E}Z_X=\mathbb{E}Z_Y=0$, $\operatorname{Var}(Z_X)=\operatorname{Var}(Z_Y)=1$, and $\operatorname{Cov}(Z_X,Z_Y)=\rho$.
\item $-1\leq\rho\leq1$, equivalently $|\rho|\leq1$.
\item $\rho=\pm1$ iff $Z_Y=\pm Z_X$ a.s., equivalently $Y=\mu_Y\pm(\sigma_Y/\sigma_X)(X-\mu_X)$ a.s., with the same sign.
\end{enumerate}
\begin{proof}
\begin{enumerate}[label=(\roman*)]
\setlength{\itemsep}{0pt}
\item $\mathbb{E}Z_X=0$, $\operatorname{Var}(Z_X)=\sigma_X^{-2}\operatorname{Var}(X)=1$, and similarly for $Z_Y$.\\
Moreover, $\operatorname{Cov}(Z_X,Z_Y)=\sigma_X^{-1}\sigma_Y^{-1}\operatorname{Cov}(X,Y)=\rho$.
\item Since $0\leq\operatorname{Var}(Z_X\pm Z_Y)=2\pm2\rho$, both $1+\rho\geq0$ and $1-\rho\geq0$, hence $|\rho|\leq1$.
\item[$(iii)$]($\Rightarrow$) If $\rho=1$, then $\mathbb{E}(Z_X-Z_Y)^2=0$, so $Z_X=Z_Y$ a.s..\\
If $\rho=-1$, then $\mathbb{E}(Z_X+Z_Y)^2=0$, so $Z_X=-Z_Y$ a.s.
\item[$(iii)$]($\Leftarrow$) If $Z_Y=\pm Z_X$, then $\rho=\operatorname{Cov}(Z_X,\pm Z_X)=\pm1$, and $(Y-\mu_Y)/\sigma_Y=\pm(X-\mu_X)/\sigma_X$\\
iff $Y=\mu_Y\pm(\sigma_Y/\sigma_X)(X-\mu_X)$ a.s.
\end{enumerate}
\hlt{Structural Commentary}: Correlation is centred covariance after unit-variance normalization.\\
\hlt{Key Insight}: Perfect correlation means no residual scatter around an affine line.
\end{proof}
\end{theorem}

\begin{example} \label{ex:covariance-correlation-bivariate-uniform-family} \hlt{Bivariate Uniform Association}\\
For $(X,Y)$ with uniform marginals on $[0,1]$, $\mathbb{E}X=\mathbb{E}Y=1/2$ and $\operatorname{Var}(X)=\operatorname{Var}(Y)=1/12$.\\
Hence $\operatorname{Cov}(X,Y)=\mathbb{E}(XY)-1/4$ and $\rho=12\operatorname{Cov}(X,Y)$.\\
If $\mathbb{E}(XY)=2/9$, then $\operatorname{Cov}(X,Y)=-1/36$, $\rho=-1/3$; if $\mathbb{E}(XY)=5/18$, then $\operatorname{Cov}(X,Y)=1/36$, $\rho=1/3$.\\
The sign of $\rho$ records negative or positive association.
\end{example}

\begin{method} \label{meth:mgf-calculus-for-classical-laws} \hlt{MGF Calculus for Classical Laws}\\
If $\psi_X(t)=\mathbb{E}e^{tX}$ exists near $0$, then $\mathbb{E}X=\psi_X'(0)$ and $\operatorname{Var}(X)=\psi_X''(0)-\psi_X'(0)^2$.
\begin{enumerate}[label=(\roman*)]
\setlength{\itemsep}{0pt}
\item If $X\sim\operatorname{Poisson}(\lambda)$, then $\psi_X(t)=\exp(\lambda(e^t-1))$, $\mathbb{E}X=\lambda$, and $\operatorname{Var}(X)=\lambda$.
\item If $X\sim\operatorname{Gamma}(\alpha,\lambda)$ with rate $\lambda$, then $\psi_X(t)=(\lambda/(\lambda-t))^\alpha$ for $t<\lambda$, $\mathbb{E}X=\alpha/\lambda$, and $\operatorname{Var}(X)=\alpha/\lambda^2$.
\item If $Z\sim N(0,1)$, then $\psi_Z(t)=e^{t^2/2}$; for $Y=\mu+\sigma Z$, $\psi_Y(t)=\exp(\mu t+\sigma^2t^2/2)$.
\end{enumerate}
\end{method}

\begin{corollary} \label{cor:mgf-independent-sum-families} \hlt{Independent Sums from MGFs}\\
Let $X\perp\!\!\!\perp Y$, and interpret all MGFs on a common neighborhood of $0$.
\begin{enumerate}[label=(\roman*)]
\setlength{\itemsep}{0pt}
\item If $\lambda,\mu>0$, $X\sim\operatorname{Poisson}(\lambda)$, and $Y\sim\operatorname{Poisson}(\mu)$, then $X+Y\sim\operatorname{Poisson}(\lambda+\mu)$.
\item If $X\sim\operatorname{Gamma}(\alpha_1,\lambda)$ and $Y\sim\operatorname{Gamma}(\alpha_2,\lambda)$ with $\alpha_1,\alpha_2,\lambda>0$, then $X+Y\sim\operatorname{Gamma}(\alpha_1+\alpha_2,\lambda)$.
\item If $\mu,\nu\in\mathbb{R}$, $\sigma,\tau>0$, $X\sim N(\mu,\sigma^2)$, and $Y\sim N(\nu,\tau^2)$, then $X+Y\sim N(\mu+\nu,\sigma^2+\tau^2)$.
\end{enumerate}
\begin{proof}
\begin{enumerate}[label=(\roman*)]
\setlength{\itemsep}{0pt}
\item $\psi_{X+Y}(t)=\psi_X(t)\psi_Y(t)=\exp(\lambda(e^t-1))\exp(\mu(e^t-1))=\exp((\lambda+\mu)(e^t-1))$, QED.
\item For $t<\lambda$, $\psi_{X+Y}(t)=\psi_X(t)\psi_Y(t)=(\lambda/(\lambda-t))^{\alpha_1+\alpha_2}$, QED.
\item $\psi_{X+Y}(t)=\psi_X(t)\psi_Y(t)=\exp(\mu t+\sigma^2t^2/2)\exp(\nu t+\tau^2t^2/2)=\exp((\mu+\nu)t+(\sigma^2+\tau^2)t^2/2)$, QED.
\end{enumerate}
\hlt{Structural Commentary}: Independence turns convolution into multiplication of MGFs.\\
\hlt{Key Insight}: Closed distribution families are stable when their MGF parameters add.
\end{proof}
\end{corollary}

\begin{method} \label{meth:moment-propagation-error-propagation} \hlt{Moment Propagation by Taylor Expansion}\\
Let the relevant variables have finite second moments, and let $g$ be $C^2$ near the relevant mean point.
\begin{enumerate}[label=(\roman*)]
\setlength{\itemsep}{0pt}
\item If $Y=g(X)$, $\mu=\mathbb{E}X$, and $\sigma^2=\operatorname{Var}(X)$, then $\mathbb{E}Y\approx g(\mu)$ and $\operatorname{Var}(Y)\approx[g'(\mu)]^2\sigma^2$.
\item The second-order mean correction is $\mathbb{E}g(X)\approx g(\mu)+\frac12g''(\mu)\sigma^2$.
\item For $Z=g(\mathbf{X})$, $\mu=\mathbb{E}\mathbf{X}$, and $\Sigma=\operatorname{Cov}(\mathbf{X})$, $\operatorname{Var}(Z)\approx\nabla g(\mu)^\top\Sigma\nabla g(\mu)$ and $\mathbb{E}Z\approx g(\mu)+\frac12\operatorname{tr}(\Sigma H_g(\mu))$.
\end{enumerate}
\hlt{Structural Commentary}: Taylor expansion propagates moments through local derivatives.\\
\hlt{Key Insight}: Smooth nonlinear functions are locally linear for variance and locally quadratic for mean correction.
\end{method}

\begin{example} \label{ex:ratio-moment-propagation} \hlt{Approximate Moments of a Ratio}\\
Let $Z=Y/X$, $\mu_X\neq0$, $\mu=(\mu_X,\mu_Y)$, and write $\sigma_X^2=\operatorname{Var}(X)$, $\sigma_Y^2=\operatorname{Var}(Y)$, $\sigma_{XY}=\operatorname{Cov}(X,Y)$.\\
For $g(x,y)=y/x$, Taylor propagation gives $\mathbb{E}Z\approx\mu_Y/\mu_X+\sigma_X^2\mu_Y/\mu_X^3-\sigma_{XY}/\mu_X^2$\\
and $\operatorname{Var}(Z)\approx\sigma_X^2\mu_Y^2/\mu_X^4+\sigma_Y^2/\mu_X^2-2\sigma_{XY}\mu_Y/\mu_X^3$. The approximation is ill-conditioned as $|\mu_X|\downarrow0$.
\end{example}
